\chapter{Applications of Phase Estimation}
% Solutions to the exercises of \cref{ch:lec05}.

\begin{solution}[Exercise~\ref{ex:lec05-1}]
On an exact eigenstate with $t$ ancilla qubits, QPE returns the best $t$-bit
dyadic approximation of the eigenphase, so the eigenphase---and via
$U=\ee^{\ii 2\pi H}$ the (rescaled) energy $\lambda$, which equals the
eigenphase---is resolved to additive error $2^{-t}$. To reach $10^{-3}$ we need
$2^{-t}\le10^{-3}$, i.e.\ $t\ge\log_2 10^3=3\log_2 10\approx9.97$, so $t=10$
ancilla qubits. If the physical spectrum spans a width $W=E_{\max}-E_{\min}$, the
affine rescaling that maps it into the phase window contracts energies by $1/W$
(up to the safety factor keeping phases in $[0,1)$); a rescaled resolution
$2^{-t}$ therefore corresponds to a physical resolution $W\cdot 2^{-t}$, so
resolving to physical accuracy $\eta$ needs $t\ge\log_2(W/\eta)$ qubits.
\end{solution}

\begin{solution}[Exercise~\ref{ex:lec05-2}]
By \cref{thm:lec05-aecost}, when $P(1-P)=\Theta(1)$ one takes
$\epsilon'=\bigo(\epsilon)$ and phase estimation of $G$ uses
$\bigo(1/\epsilon')=\bigo(1/\epsilon)$ applications of the Grover iterate. For
$\epsilon=10^{-3}$ this is $\Theta(10^{3})$ applications of $G$ (each of which
makes one query to the reflections built from $U_\phi$). A classical Monte Carlo
estimator of the Bernoulli mean $P$ has variance $P(1-P)/N=\Theta(1/N)$ after $N$
samples, so reaching additive error $\epsilon$ needs $N=\bigo(1/\epsilon^2)
=\Theta(10^{6})$ samples. The ratio of classical to quantum cost is
\[
  \frac{\bigo(\epsilon^{-2})}{\bigo(\epsilon^{-1})}=\bigo(\epsilon^{-1})
  =\Theta(10^{3}),
\]
the hallmark quadratic (here thousand-fold) speedup.
\end{solution}

\begin{solution}[Exercise~\ref{ex:lec05-3}]
A single run fails to return $\lambda_0$ with probability $1-p_0$, and the runs are
independent, so the probability that all $M$ runs miss $\lambda_0$ is
$(1-p_0)^M$. We want the complementary event to have probability $\ge\tfrac12$:
\[
  1-(1-p_0)^M\ge\tfrac12
  \iff (1-p_0)^M\le\tfrac12
  \iff M\log(1-p_0)\le-\log2
  \iff M\ge\frac{\log 2}{-\log(1-p_0)}
        =\frac{\log 2}{\log\bigl((1-p_0)^{-1}\bigr)}.
\]
For small $p_0$, $-\log(1-p_0)=p_0+\tfrac12 p_0^2+\cdots\approx p_0$, so
$M\gtrsim\log 2/p_0$. Hence $M=\lceil\log 2/p_0\rceil=\bigo(p_0^{-1})$ runs
suffice; more generally to fail with probability at most $\delta$ one takes
$M=\lceil\log(\delta^{-1})/p_0\rceil$.
\end{solution}

\begin{solution}[Exercise~\ref{ex:lec05-4}]
Let $R_\theta=\left(\begin{smallmatrix}\cos\theta&\sin\theta\\-\sin\theta&
\cos\theta\end{smallmatrix}\right)$. Its characteristic polynomial is
$(\cos\theta-\mu)^2+\sin^2\theta=0$, giving
$\mu=\cos\theta\pm\ii\sin\theta=\ee^{\pm\ii\theta}$. For $\mu=\ee^{\ii\theta}$ the
equation $(\cos\theta-\ee^{\ii\theta})v_1+\sin\theta\,v_2=0$ reads
$-\ii\sin\theta\,v_1+\sin\theta\,v_2=0$, i.e.\ $v_2=\ii v_1$, so an eigenvector is
$\tfrac{1}{\sqrt2}(1,\ii)^{\mathsf T}$; likewise $\ee^{-\ii\theta}$ has eigenvector
$\tfrac{1}{\sqrt2}(1,-\ii)^{\mathsf T}$. In ket form
$\ket{\phi_\pm}=\tfrac1{\sqrt2}(\ket{\phi_{\mathrm{good}}}\pm\ii
\ket{\phi_{\mathrm{bad}}})$. With
$\ket{\phi_0}=\sin\tfrac{\theta}{2}\ket{\phi_{\mathrm{good}}}
+\cos\tfrac{\theta}{2}\ket{\phi_{\mathrm{bad}}}$,
\[
  \braket{\phi_\pm|\phi_0}
   =\tfrac{1}{\sqrt2}\bigl(\sin\tfrac{\theta}{2}\mp\ii\cos\tfrac{\theta}{2}\bigr),
  \qquad
  \abs{\braket{\phi_\pm|\phi_0}}^2
   =\tfrac12\bigl(\sin^2\tfrac{\theta}{2}+\cos^2\tfrac{\theta}{2}\bigr)=\tfrac12,
\]
confirming the equal overlaps.
\end{solution}

\begin{solution}[Exercise~\ref{ex:lec05-5}]
A degree-$D$ polynomial $F$ with $\abs{F}\le1$ on $[-1,1]$ has a bounded
derivative there: by the Markov brothers' inequality $\abs{F'(x)}\le D^2$ on
$[-1,1]$, and more relevantly the ``transition region'' over which such a
polynomial can drop from $\approx1$ to $\approx0$ has width $\Omega(1/D)$ (the
Chebyshev polynomial $T_D$ oscillates on a scale $\sim1/D^2$ near the endpoints but
$\sim1/D$ in the bulk, and any sharper transition would force values outside
$[-1,1]$). To realize the filter we need the transition to fit inside the forbidden
window of width $\Delta$, so $1/D=\bigo(\Delta)$, i.e.\ $D=\Omega(1/\Delta)$; the
$\log(1/\epsilon')$ factor comes from pushing the approximation error down. Since
QSVT uses $\bigo(D)$ block-encoding queries, the query depth of
\cref{thm:lec05-gsp} scales as $\bigo(\Delta^{-1}\log(1/\epsilon'))$ and diverges
as $\Delta\to0$. This matches the gap lower bound $\Omega(1/\Delta)$ of
\cref{thm:lec05-lbgap}: no filter-based (or any) method can prepare the ground
state with sublinear dependence on the inverse gap.
\end{solution}

\begin{solution}[Exercise~\ref{ex:lec05-6}]
\emph{Plain QPE, $\bigo(p_0^{-2})$ (\cref{thm:lec05-gsee}, \cref{rem:lec05-overlap}):}
the $p_0^{-1}$ from the number of repetitions $M$ needed to see the ground state
at least once (each run hits $\ket{\psi_0}$ only with probability $p_0$); the
second $p_0^{-1}$ from the circuit depth $T=\bigo(\delta'^{-1})=\bigo(p_0^{-1})$
forced on us to suppress the Fej\'er-tail underestimates once $M=\bigo(p_0^{-1})$.
Their product is $TM=\bigo(p_0^{-2})$.
\emph{QSVT preparation, $\bigo(p_0^{-1/2})$ (\cref{thm:lec05-gsp}):} the sharp
filter marks the $\ket{\psi_0}$ component cleanly, so instead of \emph{sampling}
until the flag succeeds (which would cost $\bigo(p_0^{-1})$), one \emph{amplitude
amplifies} the success amplitude $\sqrt{p_0}$, needing only $\bigo(p_0^{-1/2})$
rounds (\cref{rem:lec05-sqrt}). The improvement is precisely the Grover/amplitude
amplification quadratic gain of \cref{sec:lec05-2} applied to the marked ``good''
component.
\end{solution}

\begin{solution}[Exercise~\ref{ex:lec05-7}]
\emph{Order.} $2^{1}=2$, $2^{2}=4$, $2^{3}=8$, $2^{4}=16$, $2^{5}=32\equiv11$,
$2^{6}=64\equiv1\pmod{21}$; the earlier powers are all $\ne1$, so the order is
$r=6$.

\emph{Continued fraction of $5/32$.} Flip and split:
\[
  \frac{5}{32}=\cfrac{1}{\cfrac{32}{5}}
   =\cfrac{1}{6+\cfrac{2}{5}}
   =\cfrac{1}{6+\cfrac{1}{\tfrac{5}{2}}}
   =\cfrac{1}{6+\cfrac{1}{2+\cfrac{1}{2}}},
\]
so $5/32=[0;6,2,2]$. Its convergents are
\[
  \frac{1}{6},\qquad
  \frac{1}{6+\tfrac12}=\frac{2}{13},\qquad
  \frac{5}{32}.
\]
The denominator $13$ fails ($2^{13}\bmod21=2^{6}\cdot2^{6}\cdot2\equiv1\cdot1\cdot2=2\ne1$),
and $32$ is just the ancilla resolution, not an order. The first convergent
$1/6$ has denominator $6$, and indeed $2^{6}\equiv1\pmod{21}$ with no smaller
exponent working, so $r=6$. (Equivalently, $5/32\approx0.15625$ is the closest
$5$-bit dyadic to $s/6$ with $s=1$, namely $1/6\approx0.16667$.)

\emph{Finishing Shor.} $r=6$ is even and
$2^{r/2}=2^{3}=8\not\equiv-1\equiv20\pmod{21}$, so
\[
  \gcd\bigl(2^{3}-1,\,21\bigr)=\gcd(7,21)=7,\qquad
  \gcd\bigl(2^{3}+1,\,21\bigr)=\gcd(9,21)=3,
\]
giving the factorization $21=3\times7$.
\end{solution}

\begin{solution}[Exercise~\ref{ex:lec05-8}]
The cycle generated by $a=7$ modulo $15$ is
\[
  7^{0}=1,\quad 7^{1}=7,\quad 7^{2}\equiv4,\quad 7^{3}\equiv13,
\]
returning to $7^{4}\equiv1$, so $r=4$. By \cref{lem:lec05-Ua},
\[
  \ket{u_s}=\tfrac12\Bigl(\ket{1}+\ee^{-2\pi\ii s/4}\ket{7}
              +\ee^{-2\pi\ii\cdot 2s/4}\ket{4}
              +\ee^{-2\pi\ii\cdot 3s/4}\ket{13}\Bigr),
  \qquad s=0,1,2,3,
\]
with the phase factors $\ee^{-2\pi\ii s k/4}=(-\ii)^{sk}$. Explicitly,
\[
  \ket{u_0}=\tfrac12(\ket1+\ket7+\ket4+\ket{13}),\quad
  \ket{u_1}=\tfrac12(\ket1-\ii\ket7-\ket4+\ii\ket{13}),
\]
and $\ket{u_2},\ket{u_3}$ analogously. Since $U_7$ sends
$\ket1\!\to\!\ket7\!\to\!\ket4\!\to\!\ket{13}\!\to\!\ket1$,
\[
  U_7\ket{u_1}
  =\tfrac12(\ket7-\ii\ket4-\ket{13}+\ii\ket1)
  =\ii\cdot\tfrac12(\ket1-\ii\ket7-\ket4+\ii\ket{13})
  =\ii\ket{u_1}=\ee^{2\pi\ii/4}\ket{u_1},
\]
as claimed. Finally, summing over $s$ collapses the phases by
$\tfrac14\sum_{s}\ee^{-2\pi\ii sk/4}=\delta_{k,0}$, leaving only the $k=0$ term:
$\tfrac12\sum_{s=0}^{3}\ket{u_s}=\ket{1}$. Thus preparing $\ket1$ and running QPE
is equivalent to sampling a uniformly random eigenphase $s/4$.
\end{solution}
